\documentclass[11pt,a4paper]{article}
\usepackage[utf8]{inputenc}
\usepackage[T1]{fontenc}
\usepackage{lmodern}
\usepackage[portuguese]{babel}
\usepackage[a4paper,margin=2.5cm]{geometry}
\usepackage{xcolor}
\usepackage{graphicx}
\definecolor{TextEspresso}{HTML}{4A4238}
\definecolor{NudeTaupe}{HTML}{BFAFA6}
\definecolor{SoftBlush}{HTML}{D9C5B2}
\usepackage{titlesec}
\usepackage{enumitem}
\usepackage{listings}
\usepackage{hyperref}
\hypersetup{colorlinks=true, linkcolor=TextEspresso, urlcolor=TextEspresso}

\titleformat{\section}{\color{TextEspresso}\Large\bfseries}{\thesection}{0.5em}{}[{\color{NudeTaupe}\titlerule[0.8pt]}]
\titleformat{\subsection}{\color{TextEspresso}\large\bfseries}{\thesubsection}{0.5em}{}

\lstset{
  basicstyle=\ttfamily\small,
  breaklines=true,
  frame=single,
  rulecolor=\color{NudeTaupe},
  columns=fullflexible,
  showstringspaces=false
}

\newcommand{\guiaotitle}[2]{%
  \begin{center}
  {\Large\bfseries\color{TextEspresso} #1}\\[0.3cm]
  {\normalsize\color{NudeTaupe} #2}\\[0.2cm]
  {\color{NudeTaupe}\rule{4cm}{0.5pt}}
  \end{center}
  \vspace{0.5cm}
}

\begin{document}
\guiaotitle{Guião 04 -- Containers e Isolamento de Processos}{Infraestruturas de Computação na Nuvem, Aula 04}

\section{Objetivos}
\begin{itemize}[itemsep=0.2cm]
  \item Arrancar um serviço com Docker Compose e observar o isolamento dado pelos namespaces e pelos cgroups.
  \item Comparar o mesmo processo visto de dentro e de fora do container.
  \item Verificar que um container é um processo do host, e não uma máquina virtual.
\end{itemize}

\section{Pré-requisitos e Material}
\begin{itemize}[itemsep=0.2cm]
  \item Docker e Docker Compose instalados (na VM \texttt{icn-lab} da Aula 03 ou no portátil).
  \item Conhecimentos da Aula 04: namespaces, cgroups e capabilities.
\end{itemize}

\section{Parte 1 -- Arrancar o serviço}
\begin{enumerate}[label=\arabic*.]
  \item Criar o ficheiro \texttt{docker-compose.yml} com um Redis, um volume para os dados e um limite de memória:
\begin{lstlisting}
services:
  redis:
    image: redis:7-alpine
    container_name: redis
    mem_limit: 256m
    volumes:
      - redis-data:/data
volumes:
  redis-data:
\end{lstlisting}
  \item Arrancar e confirmar o estado:
\begin{lstlisting}
docker compose up -d --build
docker compose ps
\end{lstlisting}
\end{enumerate}

\section{Parte 2 -- Explorar o isolamento por dentro}
Abrir uma shell no container e executar os comandos abaixo.
\begin{lstlisting}
docker compose exec redis sh
hostname                       # ID do container, nao o nome do portatil
ps aux                         # so o redis-server e a shell
ip addr                        # IP proprio, numa rede propria
cat /sys/fs/cgroup/memory.max  # limite imposto pelo cgroup
ls /data                       # o volume montado
redis-cli ping                 # PONG
\end{lstlisting}

\textbf{Questões.}
\begin{enumerate}[label=\arabic*.]
  \item Que namespace explica cada uma das três primeiras observações (\texttt{hostname}, \texttt{ps}, \texttt{ip addr})?
  \item Que valor tem \texttt{memory.max} e de onde vem?
  \item O que acontece aos ficheiros de \texttt{/data} se o container for removido e recriado?
\end{enumerate}

\section{Parte 3 -- O mesmo processo visto do host}
Com a shell da Parte 2 ainda aberta, noutro terminal:
\begin{lstlisting}
docker ps
docker top redis
\end{lstlisting}

\begin{enumerate}[label=\arabic*.]
  \item Registar o PID, o PPID e o utilizador do \texttt{sh} dentro do container (\texttt{ps aux}) e fora dele (\texttt{docker top}).
  \item Preencher a tabela seguinte e justificar cada linha.
\end{enumerate}

\begin{center}
\begin{tabular}{l|c|c}
  & Dentro & Fora \\ \hline
  PID & & \\
  PPID & & \\
  Utilizador & & \\
\end{tabular}
\end{center}

\textbf{Resultado esperado.}
Dentro do container o \texttt{sh} tem o PID 1, porque é o primeiro processo do namespace PID.
Fora, é o mesmo processo com um PID alto (por exemplo 48213), cujo pai é normalmente o \texttt{containerd-shim}.
O utilizador é \texttt{root} nos dois lados: por omissão o Docker não usa o namespace USER, pelo que o root do container é o root do host com menos capabilities.
É por isso que se recomenda a instrução \texttt{USER} no Dockerfile ou o modo \emph{rootless}.

Terminar o processo do lado do host (\texttt{docker kill redis}) fecha de imediato a shell do primeiro terminal, o que confirma que se trata do mesmo processo.

\textbf{Nota.}
No Docker Desktop (Mac ou Windows) os PIDs e o root mostrados pelo \texttt{docker top} são os da VM Linux onde o Docker corre, e não os do portátil.

\section{Limpeza}
\begin{lstlisting}
docker compose down -v
\end{lstlisting}

\end{document}
